\section{Discussion}
  \label{sec:discussion}

  \subsection{Status of the structural hypotheses}
    \label{sec:disc:hypotheses}

    \textbf{Hypothesis~\ref{hyp:closure}} identifies $c_\chi(n)$ with the
    isoperimetric capacity $h(G(n))$.
    It is the foundational physical postulate and currently lacks a derivation
    from the Cosmochrony axioms.

    \textbf{Hypothesis~\ref{hyp:saturation}} ($h^2 \asymp \lambda_2$) is a
    property of the graph family, consistent with but not implied by the Ramanujan
    property.
    It promotes $\Lambda_{\mathrm{proj}} \asymp h^2$ to the linear law
    $\Lambda_{\mathrm{proj}} \asymp \lambda_2$.

    \textbf{Hypothesis~\ref{hyp:saturation-rep}} ($M(\lambda_i) = k\,\dim\rho_i$,
    $k$ universal) is the only saturation condition internal to the framework that
    avoids free parameters.
    The constant $k$ cancels in all mass ratios.

    Of the three hypotheses, Hypothesis~\ref{hyp:closure} is the most in need of
    further justification.

  \subsection{Role of the parameter $p$}
    \label{sec:disc:p}

    The prime $p$ controls the Kesten--McKay support width via
    $\lambda_\pm = 1 \pm 2\sqrt{p}/(p+1)$.
    In the limit $p \to \infty$, the support collapses to $\{\lambda = 1\}$, so
    the spectral distribution becomes increasingly concentrated around the
    midpoint of the spectrum.
    Varying $p$ within the admissible range changes $c_i(p)$ quantitatively but
    not the qualitative conclusions: mass ratios remain of order one and the
    ordering inversion persists.
    The parameter $p$ should be regarded as a structural parameter of the
    relaxation-graph model.
    A natural question is whether $p$ can be selected by the framework:
    the degree $d = p+1$ is related to the number of generators of
    $\mathrm{PSL}(2,\mathbb{F}_q)$ via the four-square decomposition of
    $p$~\cite{LPS1988}; a connection to the relational valence of the Cosmochrony
    substrate would constitute a selection principle.

  \subsection{Relation with SpectralStratigraphy}
    \label{sec:disc:stratigraphy}

    The two results of~\cite{SpectralStratigraphy} used as inputs -- no-go for
    scalar calibration and three-level ADE structure -- are unaffected by the
    present analysis.

    The present work identifies two independent components of
    $n_{\mathrm{proj}}(\lambda)$: the admissibility support condition
    $\lambda \leq \Lambda_{\mathrm{proj}}$ and the saturation condition
    $N(\lambda; n) \geq k\,\dim\rho_\lambda$.
    The two mechanisms operate in different regimes.
    SpectralStratigraphy describes the regime in which $\Lambda_{\mathrm{proj}}(n)$
    evolves along the cascade; the present analysis shows that in the LPS model
    at fixed $p$ the threshold becomes asymptotically static, leaving saturation
    as the only active ordering mechanism.
    The admissibility ordering is recovered only when $\Lambda_{\mathrm{proj}}(n)$
    varies significantly, which requires a regime where $\lambda_2(n)$ is not
    asymptotically constant.

  \subsection{Interpretation of the numerical results}
    \label{sec:disc:numerical}

    The structural identity $c_2 = F_{\mathrm{KM}}(1) = 1/2$ holds for all $p$
    and all three-level ADE cases.
    It follows from the symmetry $\rho_{\mathrm{KM}}(\lambda) =
\rho_{\mathrm{KM}}(2 - \lambda)$ of the Kesten--McKay density combined with
    the fact that the central ADE level always satisfies
    $\lambda_2^{\mathrm{norm}} = 1$.
    This coincidence is not accidental but follows from the combinatorial
    normalisation of the Laplacian together with the symmetry of the Kesten--McKay
    density.
    The identity $c_2 = 1/2$ is therefore a structural alignment between the
    spectral geometry of the relaxation substrate and the representation-theoretic
    structure of the ADE groups.

  \subsection{Open directions}
    \label{sec:disc:open}

    \medskip
    \noindent\textbf{(O1) Dynamical threshold via joint $p(n)$, $q$ variation.}
    If $p$ grows with $n$, the Kesten--McKay support narrows and the ADE levels
    eventually leave the support in a rank-dependent order, restoring a non-trivial
    $\Lambda_{\mathrm{proj}}(n)$.
    Such a joint scaling would correspond physically to an increase of the
    effective relational valence of the substrate along the cascade.
    This direction is the primary candidate for addressing both the ordering
    inversion and the amplitude problem simultaneously.

    \medskip
    \noindent\textbf{(O2) Hierarchical cascade.}
    If the saturation process operates recursively across $L$ nested levels, the
    per-level ratio $r$ yields $r^L$.
    With $r \approx 0.1$ and $L = 3$, one obtains $r^3 \approx 10^{-3}$,
    approaching the lepton mass ratio $m_e/m_\tau \approx 3 \times 10^{-4}$.

    \medskip
    \noindent\textbf{(O3) Sub-linear energy-rank relation.}
    Replacing $E_n = E_{\mathrm{P}}/n$ by $E_n = E_{\mathrm{P}}/n^\alpha$ with
    $\alpha < 1$ modifies the ratios to
    $(c_i\,\dim\rho_j)^\alpha / (c_j\,\dim\rho_i)^\alpha$.
    This reduces the amplitude problem but preserves the ordering inversion, so it
    does not constitute a complete solution.
